\chapter{Introduction}
In deep learning, a key component is the training set from which our model can learn~\cite{LeCun2015}. 
It consists of the input data and the associated label. For example, if we want our model to recognise pictures of dogs and cats, the input data would be the pictures, and the associated labels would be the classification as dog or cat.
Usually, the amount of training data is required to be quite large to keep the deviation between the expected output of our model and the actual output small. 
However, the process of labelling data is not free. Often the computational cost is many times greater than the cost of generating the input itself~\cite{settles2009active}. 
A common technique to reduce the number of training samples required while still achieving high accuracy is \textit{active learning}.
Rather than just randomly labelling each piece of input data and using it as training data, 
we instead apply an algorithm that selects the samples that give us a 
high amount of information. 
This has been applied to various fields, e.\,g.\xspace drug discovery ~\cite{REKER201973} 
or text categorisation ~\cite{lewis1995sequential}. The field we are intrested in for this work is active learning for neural PDE solvers.

The solution to a PDE is typically approximated using a numerical solver, which is both resource-intensive and time-consuming. 
To reduce the effort, neural PDE solvers have been introduced~\cite{NEURIPS2023_d529b943}, 
which require a numerical solver for their training data and pass 
these results on to the neural network. Here we still have 
the problem of the high cost of labelling by numerical solvers. 
Therefore, pool-based active learning for neural PDE solvers was introduced 
by Musekamp et al. ~\cite{musekamp2024active} and applied to different neural surrogate models.

pool-based methods are one of the most common strategies for 
active learning because they give good results and are easy to implement.
However, the creation of a pool from the input space has its drawbacks. 
In particular, since the pool size is determined by a hyperparameter, 
the challenge is to find the correct pool size. If the pool is too small, 
the creation of the pool is fast, but important regions with high sample 
information from the input space will be omitted. If the pool is too large, 
it is too expensive to build. 

Instead, we can sample directly from the input space, known as query synthesis~\cite{Angluin1988}. 
The main advantage is that we can exploit the gradient of uncertainty by sampling from the largest 
increase. Unlike pool-based methods, which can become trapped in local optima, query synthesis 
allows a more global search for informative samples. This is also the main goal of this thesis: 
to replace the old pool-based method with a query synthesis strategy using the uncertainty of the model as measurement for sample information. More specifically, we want to utilize gradient information of the uncertainty, 
which should be more efficient, reducing the deviation between actual and expected results 
compared to the pool-based method, while improving the computational speed.

We begin with a brief introduction to the various fundamentals necessary to understand the 
current implementation. 
Next, we describe the methodology, followed by an explanation of the experiments and their implementation. 
We then discuss the results. Finally, we conclude with a summary of our findings and discuss 
possible directions for future research.
